%! AP Statistics — Unit 4, Lesson 4.12 — Warm-Up
\documentclass[10pt]{article}
\usepackage{apstats-article}
\usepackage{apstats-boxes}

\begin{document}

\pageheader{Unit 4, Lesson 4.12}{Warm-Up}
\namedateperiod

\vspace{0.3in}

% ── SECTION 1: Spiral Review ──────────────────────────────────────────────────
{\large\bfseries\color{navy} Part 1 \quad Spiral Review}

\vspace{0.12in}

\begin{enumerate}

    \item Let $X \sim B(8, 0.25)$.
    \begin{enumerate}[label=(\alph*), itemsep=8pt]
        \item Compute $\mu_X$ and $\sigma_X$. Show formula setup.\\
        \writeline\\[1pt]\writeline
        \item Compute $P(X = 2)$. Show the formula setup.\\
        \writeline\\[1pt]\writeline
        \item Compute $P(X \ge 3)$ using the complement rule and \texttt{binomcdf}.\\
        \writeline\\[1pt]\writeline
    \end{enumerate}

    \vspace{0.18in}

    \item For each description, circle \textbf{Binomial} or \textbf{Geometric} (or
    \textbf{Neither}). Briefly explain your choice.
    \begin{enumerate}[label=(\alph*), itemsep=10pt]
        \item A fair die is rolled until a 6 appears. $X =$ number of rolls.\\
              \textbf{Binomial} \quad / \quad \textbf{Geometric} \quad / \quad \textbf{Neither}\\
              Reason: \writeline
        \item A fair die is rolled exactly 10 times. $X =$ number of 6s.\\
              \textbf{Binomial} \quad / \quad \textbf{Geometric} \quad / \quad \textbf{Neither}\\
              Reason: \writeline
        \item Cards are drawn until an ace appears. $X =$ number of cards drawn
              (without replacement, small deck).\\
              \textbf{Binomial} \quad / \quad \textbf{Geometric} \quad / \quad \textbf{Neither}\\
              Reason: \writeline
    \end{enumerate}

    \vspace{0.18in}

    \item A basketball player makes 70\% of free throws. She keeps shooting until
    she makes her first basket.
    \begin{enumerate}[label=(\alph*), itemsep=6pt]
        \item List all outcomes where the first success is on the 3rd shot.\\
        \writeline
        \item What is the probability of each such outcome?\\
        \writeline
    \end{enumerate}

\end{enumerate}

\vspace{0.2in}

% ── SECTION 2: Looking Ahead ──────────────────────────────────────────────────
{\large\bfseries\color{navy} Part 2 \quad Looking Ahead}

\vspace{0.12in}

\noindent Problems 2 and 3 preview the key contrast for today: in a \textbf{geometric}
setting, trials continue until the first success --- $n$ is not fixed.

\vspace{0.14in}

\begin{tcolorbox}[colback=greenbg, colframe=greenacc, boxrule=0.8pt, arc=2mm,
    left=4mm, right=4mm, top=3mm, bottom=3mm]
    \small
    \begin{itemize}[leftmargin=1.3em, itemsep=8pt]
  \item Check BITS (not BINS) for a geometric setting: the T stands for
        ``trials until the first success,'' replacing N (fixed number).
  \item $P(X = k) = (1-p)^{k-1}p$ --- multiply $k-1$ failure probabilities,
        then one success.
  \item $\mu_X = 1/p$ --- larger $p$ means fewer expected trials.
    \end{itemize}
\end{tcolorbox}

\vspace{0.18in}

\noindent\textcolor{slate}{\small Write one question you have about today's lesson:}\\[8pt]
\writeline\\[2pt]\writeline

\end{document}
